We connected the partial identification of PRTEs in IV models to OT: formulating it as a CCOT problem over the joint laws compatible with the observed data, we showed that the multidimensional optimization reduces analytically to one-dimensional OT problems with product costs, yielding closed-form sharp bounds and bypassing the numerical optimization that existing computational approaches require. In a companion note we replace threshold crossing by a conditional quantile representation of the selection equation and obtain sharp closed-form bounds for continuous and multi-valued treatments, under a rank-overlap condition in the strictly monotonic case and an algebraic condition on the instrument-induced latent intervals in the ordered-choice case. The estimation theory developed here is for binary treatments, delivering $\sqrt{n}$-consistency and asymptotic normality with discrete instruments under the stated regularity conditions, and nonparametric convergence rates with continuous ones.

The central insight, demonstrated theoretically and empirically, is that preserving full distributional information is tractable and valuable for tight partial identification: moment-relaxation approaches reduce the observed data to conditional expectations and can produce substantially wider bounds, whereas CCOT enforces compatibility with the entire observed distribution while keeping the optimization transparent. We expect this to extend beyond the IV setting studied here.


\paragraph{Acknowledgments.} Jose Blanchet gratefully acknowledges support from the U.S. Department of Defense through the Air Force Office of Scientific Research award FA9550-20-1-0397 and ONR award 1398311, as well as support from NSF grants 2229012, 2312204, and 2403007.
Vasilis Syrgkanis gratefully acknowledges support from NSF Award IIS-2337916, an Amazon Research Award, and a Google Research Award.
